% !PRINTED 229
\weberpage{229}
In addition, we set:
\weq{4}{P_{11}\,\eta(\omega)^{\frac{n-1}{2}} = e^{\frac{\pi i}{6}\mu'(\mu+\mu'\omega)\frac{n^2-1}{n}}\;\mathop{\Pi}\limits_{1,\frac{n-1}{2}}^{h}\vartheta_{11}(2h\varpi).}

These quantities $P$ can be expressed by the roots of the division equation if, applying (3), one forms the quotients
\[
\frac{P_{00}^2}{P_{01}P_{10}},\quad \frac{P_{10}^2}{P_{01}P_{00}},\quad \frac{P_{01}^2}{P_{10}P_{00}},\quad \frac{P_{11}^3}{P_{00}P_{01}P_{10}}
\]
and by means of the formulae of § 42 introduces elliptic functions.

One obtains in this way
\weq{5}{\begin{aligned}
(-1)^{\frac{n^2-1}{8}}\,2^{\frac{n-1}{2}}P_{00}^3 &= \mathop{\Pi}\limits_{1,\frac{n-1}{2}}^{h}\frac{\operatorname{dn}^2h\Omega}{\operatorname{cn}h\Omega},\\
(-1)^{\frac{n^2-1}{8}}\,2^{\frac{n-1}{2}}P_{10}^3 &= \mathop{\Pi}\limits_{1,\frac{n-1}{2}}^{h}\frac{\operatorname{cn}^2h\Omega}{\operatorname{dn}h\Omega},\\
(-1)^{\frac{n^2-1}{8}}\,2^{\frac{n-1}{2}}P_{01}^3 &= \mathop{\Pi}\limits_{1,\frac{n-1}{2}}^{h}\frac{1}{\operatorname{cn}h\Omega\,\operatorname{dn}h\Omega},
\end{aligned}}
\weq{6}{(-1)^{\frac{n^2-1}{8}}P_{11}^3 = (\varkappa\varkappa')^{\frac{n-1}{2}}\mathop{\Pi}\limits_{1,\frac{n-1}{2}}^{h}\frac{\operatorname{sn}^3h\Omega}{\operatorname{cn}h\Omega\,\operatorname{dn}h\Omega}.}

These expressions now show that $P_{00}^3$, $P_{10}^3$, $P_{01}^3$ \gesperrt{are the roots of transformation equations (modular equations)}. $P_{11}^6$ is the root of a multiplier equation, and if $n$ is a square number, then $P_{11}^3$ too is \gesperrt{the root of a multiplier equation}.

One can combine the functions $P$ with one another in manifold ways in order to obtain new transformation equations. We adduce the following:
\weq{7}{\begin{aligned}
\frac{P_{10}}{P_{00}} &= \mathop{\Pi}\limits_{1,\frac{n-1}{2}}^{h}\frac{\operatorname{cn}h\Omega}{\operatorname{dn}h\Omega},\\
\frac{P_{01}}{P_{00}} &= \mathop{\Pi}\limits_{1,\frac{n-1}{2}}^{h}\frac{1}{\operatorname{dn}h\Omega},
\end{aligned}}
%CONT>>
